
% BEGIN bib
 
\par\medskip\noindent\textbf{References}\par
  
\par\noindent  [1]  Samson Abramsky, Anuj Dawar, and Pengming Wang.  The pebbling comonad in finite model theory.  In Proceedings of LICS, pages 1–12. IEEE Computer Society, 2017.  doi:10.1109/LICS.2017.8005129.  
 
\par\noindent  [2]  Samson Abramsky and Luca Reggio.  Arboreal categories: An axiomatic theory of resources.  Log. Methods Comput. Sci., 19(3), 2023.  doi:10.46298/LMCS-19(3:14)2023.  
 
\par\noindent  [3]  Samson Abramsky and Nihil Shah.  Relating structure and power: Comonadic semantics for computational resources.  J. Log. Comput., 31(6):1390–1428, 2021.  doi:10.1093/LOGCOM/EXAB048.  
 
\par\noindent  [4]  John C. Baez, Joe Moeller, and Todd Trimble.  2-rig extensions and the splitting principle.  arXiv:2410.05598 [math.CT], 2024.  
 
\par\noindent  [5]  Lenore Blum, Mike Shub, and Steve Smale.  On a theory of computation and complexity over the real numbers.  Bulletin of the American Mathematical Society, 21(1):1–46, 1989.  
 
\par\noindent  [6]  Martin Brandenburg.  Tensor categorical foundations of algebraic geometry.  Ph.d. thesis, University of Münster, 2014.  
 
\par\noindent  [7]  Martin Brandenburg and Alexandru Chirvasitu.  Tensor functors between categories of quasi-coherent sheaves.  Journal of Algebra, 399:675–692, 2014.  
 
\par\noindent  [8]  Aurelio Carboni, Stephen Lack, and R.F.C. Walters.  Introduction to extensive and distributive categories.  Journal of Pure and Applied Algebra, 84(2):145–198, 1993.  
 
\par\noindent  [9]  Baptiste Chanus and Damiano Mazza.  A categorical approach to describing complexity classes.  In preparation. Available on the authors’ web page, 2024.  
 
\par\noindent  [10]  Alexandru Chirvasitu and Theo Johnson-Freyd.  The fundamental pro-groupoid of an affine 2-scheme.  Appl. Categ. Structures, 21(5):469–522, 2013.  doi:10.1007/S10485-011-9275-Y.  
 
\par\noindent  [11]  Anuj Dawar, Tomás Jakl, and Luca Reggio.  Lovász-type theorems and game comonads.  In Proceedings of LICS, pages 1–13. IEEE, 2021.  doi:10.1109/LICS52264.2021.9470609.  
 
\par\noindent  [12]  Nikolai Durov.  New Approach to Arakelov Geometry.  arXiv 0704.2030 [math.AG], 2007.  
 
\par\noindent  [13]  Nikolai Durov.  Classifying vectoids and operad kinds.  Proc. Steklov Inst. Math., 273:48–63, 2011.  
 
\par\noindent  [14]  Ronald Fagin.  Generalized first-order spectra and polynomial-time recognizable sets.  In Complexity of Computation, SIAM-AMS Proccedings, volume 7, pages 43–73, 1974.  
 
\par\noindent  [15]  Erich Grädel and Klaus Meer.  Descriptive complexity theory over the real numbers.  In Proceedings of STOC, pages 315–324. ACM Press, 1995.  doi:10.1145/225058.225151.  
 
\par\noindent  [16]  Neil Immerman.  Relational queries computable in polynomial time.  Information and Control, 68(1):86–104, 1986.  doi:10.1016/S0019-9958(86)80029-8.  
 
\par\noindent  [17]  Neil Immerman.  Descriptive complexity.  Graduate texts in computer science. Springer, 1999.  doi:10.1007/978-1-4612-0539-5.  
 
\par\noindent  [18]  Bart P. F. Jacobs.  Categorical Logic and Type Theory, volume 141 of Studies in logic and the foundations of mathematics.  North-Holland, 2001.  
 
\par\noindent  [19]  Peter T. Johnstone.  Sketches of en Elephant. A Topos Theory Compendium. Volume 2.  Oxford University Press, 2002.  
 
\par\noindent  [20]  Melven R. Krom.  The decision problem for a class of first-order formulas in which all disjunctions are binary.  Zeitschrift für Mathematische Logik und Grundlagen der Mathematik, 13(1–2):15–20, 1967.  
 
\par\noindent  [21]  Saunders Mac Lane and Ieke Moerdijk.  Sheaves in Geometry and Logic.  Springer, 1994.  
 
\par\noindent  [22]  Yoàv Montacute and Glynn Winskel.  Concurrent games over relational structures: The origin of game comonads.  In Proceedings of LICS, pages 58:1–58:14. ACM, 2024.  doi:10.1145/3661814.3662075.  
  
% END bib
